\begin{tikzpicture}[
    font=\small,
    >=stealth,
    block/.style={
        draw, 
        rounded corners=12pt, 
        minimum width=2.4cm, 
        minimum height=2.4cm, 
        align=center, 
        thick
    }
]

    % ==========================================
    % 1. RAMA SUPERIOR: VENA / OPTIMIZADOR
    % ==========================================

    % Entrada 1: Vena de interés a evaluar
    \node[block] (in_vessel) at (0.65, 3.0) {};
    
    % Dibujo interno de la vena de interés (ramificación)
    \begin{scope}[shift={(in_vessel.center)}]
        \draw[line width=2.5pt, line cap=round, draw=black!80] 
            (-0.4, -0.6) .. controls (-0.1, -0.1) .. (0.0, 0.2);
        \draw[line width=2pt, line cap=round, draw=black!80] 
            (0.0, 0.2) .. controls (0.2, 0.4) .. (0.5, 0.6);
        \draw[line width=1.5pt, line cap=round, draw=black!80] 
            (0.0, 0.2) .. controls (-0.2, 0.4) .. (-0.4, 0.7);
        \draw[line width=1.2pt, line cap=round, draw=black!80] 
            (-0.15, -0.15) .. controls (0.2, 0.0) .. (0.4, 0.1);
    \end{scope}

    \node[above=0.05cm of in_vessel, font=\bfseries\footnotesize] {Vena de Interés};

    % Bloque Optimizador
    \node[block, right=1.2cm of in_vessel] (opt) {
        \textbf{Optimizador}
    };

    % Flecha Vena -> Optimizador
    \draw[thick, ->] (in_vessel.east) -- (opt.west);


    % ==========================================
    % 2. RAMA INFERIOR: SENSORES / FBPINN / EXTENSIÓN
    % ==========================================

    % Entrada 2: Geometría de Sensores y Fuente
    \node[block] (in_sens) at (0.65, -1.0) {};
    
    % Dibujo interno de los 20 sensores concentricos + pulso central
    \begin{scope}[shift={(in_sens.center)}]
        % Círculo guía punteado para los sensores
        \draw[dashed, gray!60] (0,0) circle (0.75cm);
        
        % 20 puntos concentricos (sensores)
        \foreach \angle in {0,18,...,342} {
            \fill[black!80] (\angle:0.75cm) circle (0.04cm);
        }
        
        % Objeto / pulso medio difuminado en el centro
        \node[shading=radial, inner color=red!90!black, outer color=white, circle, minimum size=0.55cm] at (0,0) {};
        \fill[red!90!black] (0,0) circle (0.06cm);
    \end{scope}

    \node[above=0.05cm of in_sens, font=\bfseries\footnotesize] {Geometría Sensores};

    % Bloque FBPINN
    \node[block, right=1.2cm of in_sens] (fbpinn) {
        \textbf{FBPINN}
    };

    % Bloque Extensión a Dominio Mayor
    \node[block, right=1.2cm of fbpinn] (ext) {
        \textbf{Extensión a}\\[0.1cm]
        \textbf{Dominio Mayor}
    };

    % Flechas Rama Inferior
    \draw[thick, ->] (in_sens.east) -- (fbpinn.west);
    \draw[thick, ->] (fbpinn.east) -- (ext.west);


    % ==========================================
    % 3. UNIFICACIÓN: FUNCIONES BASE Y SUMATORIA
    % ==========================================

    % Bloque Generación de Funciones Base
    \node[block, minimum width=2.6cm] (gen) at (11.7, 0.99) {
        \textbf{Generación de}\\[0.1cm]
        \textbf{Funciones Base}
    };

    % Flecha Optimizador -> Entrada superior de 'gen'
    \draw[thick, ->] (opt.east) -- node[above, font=\small] {$a_i, \mu_i$} (opt.east -| gen.north) -- (gen.north);

    % Flecha Extensión -> Entrada inferior de 'gen'
    \draw[thick, ->] (ext.east) -| (gen.south);

    % Bloque Sumatoria (Sigma)
    \node[block, minimum width=1.8cm, minimum height=1.8cm, right=0.87cm of gen] (sum) {
        \Huge $\sum$
    };

    % Flecha gen -> sum
    \draw[thick, ->] (gen.east) -- (sum.west);


    % ==========================================
    % 4. SALIDA FINAL: SINOGRAMA RESULTANTE
    % ==========================================

    \node[block, right=0.64cm of sum, fill=green!5!white] (out) {
        \textbf{Sinograma }
      
    };

    % Flecha sum -> Salida
    \draw[thick, ->] (sum.east) -- (out.west);

\end{tikzpicture}
